\section*{怎么读：路线与数学预备}
\phantomsection\label{front:reading-routes}
\begingroup
\setlength{\parskip}{4pt}

本书沿一条主线前进：训练参数，把数字扩展为矩阵，
让 Token 读取上下文并预测下一 Token，最后接入助手与 Agent。
首次阅读先抓住每一步的输入、输出和用途，再追踪矩阵元素。

\subsection*{三条阅读路线}
\begin{itemize}[itemsep=3pt,parsep=0pt]
\item \textbf{速读：}看每部分地图、章首问题、形状图与章末小结，
先解释训练、矩阵、Token、Attention、Transformer 和外部系统怎样相连。
\item \textbf{标准：}按正文顺序阅读，亲手算损失、梯度更新、交叉熵与注意力；
做完各章三题，用题后的一行答案提示核对，再读附录 D 的解题过程。
\item \textbf{动手：}按章末题号做附录 C，运行附录 B 的 PyTorch 实验，
用 Notebook 滑块观察参数与概率变化。联想记忆的提高专题见附录 H--K。
\end{itemize}
视觉部分用 CNN 与 ViT 连接图像和语言；助手部分区分模型计算、证据与工具；
架构部分比较组件替换的收益与代价。按页眉定位章节，以能否独立复算决定阅读速度。

\subsection*{数学预备：先核对四件事}
需要会解一元一次方程、计算负数和平方，并理解函数、对数与一元导数；
不熟悉就先读\hyperref[app:math-tools]{附录 A.2 数学速查}。

\noindent 代入 $2(3)+1=7$；平方 $(-2)^2=4$；对数 $\log(e^2)=2$。\par
\noindent 链式法则：$z=wx+b,\ L=z^2/2$ 时，$\partial L/\partial w=zx$。

现在就试一次：样本 $(1,3),(2,5),(3,7)$ 满足 $\widehat y=2x+1$，
代入 $x=4$ 得 $9$。接下来学怎样用损失自动找到参数。

\subsection*{表示约定：数学从 1 起，代码索引减 1}
\textbf{ID、位置和矩阵行列都从 1 数。}教学词表是
$1=\text{其他}$、$2=\text{春天}$、$3=\text{。}$、$4=\text{花}$、$5=\text{开了}$。
第一位置预测“花”的概率写作 $\widehat P_{1,4}$，代码是 \texttt{probs[0,3]}。

\textbf{每个 Token 占一行，投影写成 $XW$。}
Q 蓝色、K 青绿色、V 橙色，同时保留文字标签。
Python 只在查表和标签索引时减 1；代码实验见\hyperref[app:pytorch-lab]{附录 B}。

\textbf{认清例子身份。}“独立例”使用自己的参数；教学报表与奖励均为虚构数据。
\endgroup
